\chapter{KẾT LUẬN VÀ HƯỚNG PHÁT TRIỂN}\label{Chapter5}

\section*{Tóm tắt chương}

Chương này tổng kết nội dung khóa luận, nêu các đóng góp chính và đề xuất hướng phát triển tiếp theo cho hệ thống Decentralized Federated Distillation kết hợp Class-Incremental Learning.

\section{Kết luận}

Khóa luận đã trình bày một cơ chế học liên kết hỗ trợ học liên kết giữa các kiến trúc mô hình khác nhau (model heterogeneity), được được hiện thực (ký hiệu (multiple) trong bảng so sánh \ref{tab:cil_results_task6}) cho phép các client triển khai những mô hình hoàn toàn khác nhau mà vẫn cùng tham gia một quá trình học liên kết thống nhất. Xuyên suốt mọi task, hiệu quả phân loại của kịch bản đa kiến trúc mô hình liên tục tiệm cận các phương pháp đơn cấu trúc và thậm chí thuộc nhóm dẫn đầu ở một số task như task 2 với Accuracy tại task 2 ở $0.9245$, gần tương đương với FEAT ở $0.9325$, một phương pháp State-of-the-art đơn kiến trúc. Phương pháp trong kịch bản này cho thấy chỉ số Macro Precision cao thường xuyên dẫn đầu xuyên suốt 7 task học tiệm tiến. Đây là kết quả thực nghiệm khi có đến $25\%$ số client sử dụng mô hình nhỏ là CNN \ref{tab:cnn_arch}, và các mô hình không chỉ khác nhau về số đối mà còn khác nhau về mặt kiến trúc.

Ở kịch bản đơn kiến trúc, phương pháp đề xuất liên tục dẫn đầu so với các phương pháp khác, và luôn nằm ở mức tương đương so với FEAT. Điều này cho thấy hiệu quả học tiệm tiến rõ rệt của phương pháp đề xuất.

Tuy nhiên khác với FEAT, trong phương pháp đề xuất, do tri thức được trao đổi dưới dạng logits trên tập dữ liệu công khai thay vì trọng số mô hình, hệ thống đề xuất hoàn toàn tránh được các tấn công inference lên trọng số mô hình, giảm thiểu bề mặt của tấn công membership inference và hoàn toàn hỗ trợ được nhiều kiến trúc mô hình khác nhau trong quá trình học liên kết. Cùng với đó, cơ chế trao đổi dựa trên dữ liệu công khai mở ra khả năng triển khai phi tập trung mà không cần một máy chủ tổng hợp tham số trung tâm.

Thực nghiệm của hệ thống cho thấy độ chính xác cao, vượt qua nhiều phương pháp học tiệm tiến trong khi giải quyết nhiều khó khăn thực tế của các hệ thống học máy liên kết. Bằng cách ứng dụng phương pháp học tiệm tiến dựa trên tái diễn dữ liệu và kiến trúc dynamic, cũng như nhiều kỹ thuật tiên tiến có chi phí thấp như EVA và EKD, khóa luận đã thành công trình bày một hệ thống mang tính thực tiễn với độ hiệu quả cao cho bài toán học liên kết phát hiện xâm nhập.

\section{Đóng góp chính}

\begin{itemize}
    \item \textbf{Hệ thống học liên kết phân tán dựa trên học chắt lọc kết hợp học tiệm tiến theo nhãn}: Khóa luận đã xây dựng một hệ thống hoàn chỉnh cho phép mô phỏng và đánh giá các phương pháp học liên kết trong ngữ cảnh học tiệm tiến lớp, so sánh với $9$ phương pháp học tiệm tiến của cả ba loại: dựa trên ràng buộc (regularization), tái diễn dữ liệu (replay/exemplar rehearsal), và kiến trúc dynamic (dynamic architecture) cùng phương pháp đề xuất dựa trên chắt lọc liên kết. Hệ thống được hiện thực chủ yếu trên PyTorch với khả năng chuyển đổi linh hoạt giữa hai backend TensorFlow và PyTorch, hỗ trợ cả kịch bản đơn kiến trúc và đa kiến trúc, đồng thời cung cấp cơ chế checkpoint để phục hồi sau sự cố trong quá trình huấn luyện kéo dài, file log chi tiết với mốc đo thời gian.
    
    \item \textbf{Phương pháp đề xuất}: Phương pháp đề xuất đã giải quyết được nhiều khó khăn mang tính thực tiễn cao: khả năng hỗ trợ kiến trúc mô hình khác nhau, cải thiện bảo mật dữ liệu trước các tấn công suy ngược mô hình và gây khó khăn nhiều hơn cho các kiểu tấn công membership inference trong khi vẫn giữ được khả năng phân loại tiệm cận hoặc cao hơn các phương pháp học tiệm tiến State-of-the-art như FEAT.
    
    % \item \textbf{Chi phí truyền thông co giãn được}: Phương pháp đề xuất đã biến chi phí truyền thông từ một ràng buộc cứng thành một tham số thiết kế có thể tinh chỉnh. Thay vì buộc client phải tải lên toàn bộ trọng số mô hình ở mỗi vòng như các phương pháp dựa trên tham số, phương pháp đề xuất trao đổi logits trên dữ liệu công khai với lượng thông tin tỉ lệ thuận với số mẫu công khai được chọn, cho phép người triển khai chủ động đánh đổi giữa băng thông và chất lượng đồng thuận.
    
    \item \textbf{Đánh giá thực nghiệm toàn diện trên CICIoT2023}: Khóa luận đã tiến hành đánh giá thực nghiệm toàn diện trên tập dữ liệu CICIoT2023 với $34$ lớp tấn công phân bố trên $7$ task nối tiếp, trong đó số lượng client tăng dần từ $40$ lên $100$ để mô phỏng quá trình mở rộng mạng lưới IDS liên kết. Kết quả thực nghiệm đã làm rõ khả năng của từng phương pháp trong việc học tiệm tiến tiếp nhận các lớp tấn công mới mà vẫn duy trì tri thức đã học, .
\end{itemize}

\section{Hướng phát triển}

Bên cạnh các khó khăn đã được giải quyết hiệu quả bằng phương pháp đề xuất, khóa luận vẫn còn nhiều không gian phát triển để minh chứng và cải thiện tốt hơn tính ứng dụng:

\begin{itemize}
    % \item \textbf{Triển khai trên hệ thống IDS thực tế}: Hướng phát triển ưu tiên hàng đầu là đưa hệ thống từ môi trường mô phỏng vào triển khai thực tế trên các hệ thống phát hiện xâm nhập trong doanh nghiệp hoặc tổ chức. Việc này đòi hỏi giải quyết các thách thức về tích hợp với hạ tầng mạng hiện có, xử lý lưu lượng mạng thời gian thực với độ trễ thấp, và đảm bảo khả năng mở rộng khi số lượng client và lớp tấn công tăng lên. Thử nghiệm trên môi trường thực tế cũng sẽ cung cấp dữ liệu đánh giá phong phú hơn, bao gồm các kiểu tấn công mới nổi và phân phối dữ liệu phức tạp hơn so với tập dữ liệu CICIoT2023.
    
    \item \textbf{Tích hợp với hệ thống SIEM/SOC}: Hệ thống có thể được tích hợp với các hệ thống Security Information and Event Management (SIEM) hoặc Security Operations Center (SOC) để cung cấp khả năng phát hiện xâm nhập phân tán và thích nghi theo thời gian. Trong kịch bản này, mỗi tổ chức tham gia sẽ vận hành một client của hệ thống, và tri thức về các kiểu tấn công mới được chia sẻ thông qua cơ chế chắt lọc liên kết mà không làm lộ dữ liệu nhạy cảm. Việc tích hợp này cũng đòi hỏi phát triển giao diện người dùng trực quan để các nhà phân tích bảo mật có thể theo dõi hiệu năng của mô hình, điều tra các cảnh báo sai, và cập nhật chính sách phản ứng tự động.
    
    % \item \textbf{Kiểm chứng thực nghiệm khả năng phi tập trung}: Mặc dù phương pháp đề xuất đã được phân tích ở mức lý thuyết về khả năng triển khai phi tập trung, việc kiểm chứng thực nghiệm vẫn chưa được thực hiện. Hướng phát triển tiếp theo là xây dựng một kiến trúc mạng ngang hàng (peer-to-peer) cho phép các client trao đổi logits trực tiếp mà không cần máy chủ trung tâm, đồng thời đánh giá hiệu năng của hệ thống trong các điều kiện mạng không ổn định, độ trễ cao, và một số client có thể rời khỏi hoặc gia nhập hệ thống bất cứ lúc nào.
    
    \item \textbf{Gia tăng hiệu năng}: Mặc dù phương pháp đề xuất đã đạt được các tính chất hệ thống mong muốn, độ chính xác phân loại xét riêng trong việc phát hiện các nhãn tấn công vẫn còn đạt hiệu quả trung bình, khi các mô hình chỉ phát hiện được các tấn công thường thấy (như tấn công DDoS) mà lại phân loại kém hơn trên các kiểu tấn công hiếm (XSS, Dictionary Brute Forcing,...), dẫn tới chỉ số Macro F1 còn chưa tối ưu. Đây là một vấn đề sinh ra từ đặc tính mất cân bằng nhãn rất cao của tập dữ liệu CICIoT23, và giải quyết vấn đề này là một hướng phát triển hoàn toàn cần thiết của khóa luận này.
    
    % \item \textbf{Đánh giá trên nhiều tập dữ liệu IDS khác}: Để khẳng định tính tổng quát của hệ thống, cần tiến hành đánh giá trên nhiều tập dữ liệu phát hiện xâm nhập khác nhau, bao gồm các tập dữ liệu với đặc trưng mạng đa dạng hơn (ví dụ: NSL-K, UNSW-NB15, CSE-CIC-IDS2018) và các kịch bản tấn công phức tạp hơn (ví dụ: tấn công zero-day, tấn công đa giai đoạn). Việc này cũng sẽ giúp xác định các trường hợp sử dụng phù hợp nhất cho từng phương pháp học tiệm tiến.
    
    % \item \textbf{Tối ưu chi phí truyền thông và tính toán}: Mặc dù phương pháp đề xuất đã cho phép điều chỉnh linh hoạt chi phí truyền thông, việc tối ưu hóa cả chi phí truyền thông lẫn chi phí tính toán ở phía client vẫn là một hướng nghiên cứu quan trọng. Các kỹ thuật như lượng tử hóa logits, nén dữ liệu công khai, và học tăng cường để tự động điều chỉnh tỉ lệ dữ liệu công khai theo điều kiện mạng có thể được khám phá để giảm thiểu tổng chi phí vận hành.
    
    \item \textbf{Bổ sung phân tích khả giải cho dự đoán IDS}: Trong bối cảnh ứng dụng vào phát hiện xâm nhập, khả năng giải thích dự đoán của mô hình có vai trò quan trọng để các nhà phân tích bảo mật có thể hiểu được lý do đằng sau các cảnh báo và đưa ra quyết định phản ứng phù hợp. Hướng phát triển này bao gồm việc tích hợp các kỹ thuật như SHAP, LIME, hoặc attention visualization vào hệ thống, qua đó cung cấp thông tin chi tiết về các đặc trưng quan trọng nhất trong việc phân loại từng kiểu tấn công.
    
    \item \textbf{Hoàn thiện cơ chế chống tấn công poisoning}: Trong môi trường thực tế, một số client có thể bị xâm nhập hoặc hoạt động bất thường (Byzantine actors), dẫn đến việc gửi thông tin sai lệch cho hệ thống. Hướng phát triển tiếp theo là hoàn thiện tích hợp một cơ chế phát hiện và cách ly client bất thường, đồng thời thiết kế các giao thức đồng thuận bền vững trước các kiểu tấn công nhằm vào kiến trúc phân tán như Free-rider attack hay Sybil attack \cite{douceur2002sybil}
\end{itemize}
